% Formatting proof only. No measured experimental results are supplied.
% Uses the official Springer sample's 2026-09-03 font packages unchanged.
\documentclass{llncs}
\usepackage[T1]{fontenc}
\usepackage{newtxtext}
\usepackage[varvw]{newtxmath}
\usepackage{graphicx}
\usepackage[hidelinks]{hyperref}
\hypersetup{pdftitle={Anonymous Formatting Skeleton},pdfauthor={},
  pdfsubject={Template compilation proof only},pdfkeywords={},pdfcreator={LaTeX}}
\pdfinfoomitdate=1
\pdftrailerid{}
\begin{document}
\title{Decision Support for Urban Solar Access Anonymous Formatting Skeleton}
\author{}
\institute{}
\maketitle
\begin{abstract}
This document is a formatting proof for a future research manuscript. It tests
an anonymous title block, editable equations and tables, a vector diagram,
figure alternative text, numbered citations, and the bibliography style in the
current Springer proceedings class. It contains no experimental results and
must not be submitted as a research paper. The final abstract must be written
from reviewed evidence after the geometry, temporal, and matched runtime
checks have been resolved. Claims require traceable sources and a distinction
between historical evidence and new measurements. The sections below identify
their intended scope without filling missing observations with illustrative
numbers. This short proof is deliberately below the conference's minimum
character count; successful compilation establishes only that the template
can be built. Author approval, scientific review, disclosure placement, and
a submission-stage anonymity check remain necessary. The provisional
disclosure records the artificial-intelligence assistance used for this proof
and requires confirmation of the conference's placement requirements.
\keywords{Decision support \and Solar access \and Reproducibility}
\end{abstract}

\section{Introduction}
Define the decision problem, setting, and evidence-supported contribution.
The Ladybug paper describes an environmental analysis
tool~\cite{roudsari2013ladybug}. It demonstrates citation formatting here;
it does not validate this study.

\section{Methods}
\subsection{Decision Variables and Sunlight Duration}
Equation~\ref{eq:selection} illustrates the editable binary decision model.
The eligible set $E$ must be defined by explicit research constraints; it does
not represent legal permission to alter a building.
\begin{equation}
 z_i\in\{0,1\}\quad(i\in E),\qquad \sum_{i\in E}z_i=2.
 \label{eq:selection}
\end{equation}
Table~\ref{tab:symbols} defines the symbols used in
Eq.~\ref{eq:duration}. For fixed sample times, the direct-sunlight duration at
receiver $j$ is represented by
\begin{equation}
 D_j(z)=\sum_{t\in T}\Delta t\,v_{jt}(z)
 \mathbf{1}\{\alpha_t>0\}.
 \label{eq:duration}
\end{equation}
State the interval interpretation, solar-position source, visibility
tolerances, receiver mapping, and geometry treatment. The notation asserts
no agreement with an external solver.

\begin{table}[t]
\caption{Editable notation example with no experimental values}
\label{tab:symbols}
\centering
\begin{tabular}{l@{\quad}l@{\quad}l}
\hline
Symbol & Meaning & Unit \\
\hline
$z_i$ & Selection of eligible building $i$ & dimensionless \\
$T$ & Fixed set of analysis samples & index set \\
$\Delta t$ & Duration represented by each sample & h \\
$v_{jt}(z)$ & Direct visibility at receiver $j$ and sample $t$ & binary \\
$\alpha_t$ & Solar elevation at sample $t$ & angle \\
$D_j(z)$ & Direct-sunlight duration at receiver $j$ & h \\
\hline
\end{tabular}
\end{table}

Figure~\ref{fig:workflow} illustrates the evidence sequence. Describe
containment, volume, matched runtime, and validation checks separately.

\begin{figure}[t]
\centering
\setlength{\unitlength}{1mm}
\begin{picture}(114,16)
\put(0,3){\framebox(30,10){Fixed protocol}}
\put(30,8){\vector(1,0){10}}
\put(40,3){\framebox(30,10){Audits and tests}}
\put(70,8){\vector(1,0){10}}
\put(80,3){\framebox(34,10){Evidence review}}
\end{picture}
\Description{Three boxes connected from left to right: Fixed protocol, Audits and tests, and Evidence review. This is a conceptual workflow with no measured data.}
\caption{Conceptual vector workflow with no measured data. AI tool: OpenAI assistant. Prompt: show fixed protocol, audits and tests, and evidence review in sequence.}
\label{fig:workflow}
\end{figure}

\section{Results, Limitations, and Conclusions}
Insert reviewed evidence, retain unresolved effects, and separate measurements
from extrapolations. Write conclusions only after scientific review.

\section*{Provisional AI Use Disclosure}
An OpenAI assistant~\cite{openai2026} generated this proof, including its
prose, equations, table, and conceptual diagram. No experimental values were
generated. Final AI-use disclosure and its anonymous placement require author
review. Human authors must also declare competing interests; none are
asserted on their behalf here.

\bibliographystyle{splncs04}
\bibliography{references}
\end{document}

